% Part: propositional-logic
% Chapter: syntax-and-semantics
% Section: introduction

\documentclass[../../../include/open-logic-section]{subfiles}

\begin{document}

\olfileid{pl}{syn}{int}

\olsection{Inleiding}

De propositielogica bestudeert !!{formula}s die uit
!!{propositional variable}s worden opgebouwd met de propositionele
connectieven $\lnot$, $\land$, $\lor$, $\lif$ en $\liff$. Intuïtief
staat !!a{propositional variable}~$p$ voor een zin of propositie die
waar of onwaar is. Zodra de ``waarheidswaarde'' van de
!!{propositional variable}s in !!a{formula} is bepaald, ligt ook de
waarheidswaarde vast van elke !!{formula} die daaruit met propositionele
connectieven is gevormd. We noemen de propositielogica
\emph{waarheidsfunctioneel}, omdat haar semantiek door functies van
waarheidswaarden wordt gegeven. In de propositielogica laten we nadere
bepalingen van waarheid en onwaarheid dus buiten beschouwing. We vragen
bijvoorbeeld niet of iets noodzakelijk waar is in plaats van slechts
contingent waar, of bekend is dat iets waar is, of iets nu waar is in
plaats van vroeger waar te zijn geweest of later waar te zullen zijn.
We beschouwen slechts twee waarheidswaarden, waar ($\True$) en onwaar
($\False$), en sluiten daarmee uit dat een bewering noch waar noch
onwaar, of slechts half waar, zou kunnen zijn. Ook beperken we ons tot
connectieven waarvoor de waarheidswaarde van een ermee opgebouwde
!!{formula} volledig door de waarheidswaarden van haar delen wordt
bepaald (en bijvoorbeeld niet door haar betekenis). Of de
waarheidswaarde van voorwaardelijke zinnen in het Engels in deze
zin waarheidsfunctioneel is, is omstreden. De materiële implicatie
$\lif$ is dat wel; andere logica's behandelen voorwaardelijke
verbindingen die niet waarheidsfunctioneel zijn.

Om de theorie en metatheorie van de waarheidsfunctionele
propositielogica te ontwikkelen, moeten we eerst de syntaxis en
semantiek van haar uitdrukkingen definiëren. We beschrijven één manier
om met connectieven !!{formula}s uit !!{propositional variable}s op te
bouwen. Andere definities zijn mogelijk. Andere systemen kiezen andere
symbolen, nemen andere verzamelingen connectieven als primitief en
gebruiken haakjes anders (of helemaal niet, zoals in de zogenoemde
Poolse notatie). Alle benaderingen hebben echter gemeen dat de
vormingsregels de verzameling !!{formula}s \emph{inductief} definiëren.
Wanneer dat zorgvuldig gebeurt, kan elke uitdrukking volgens de
vormingsregels in wezen maar op één manier zijn ontstaan. De inductieve
definitie levert dan \emph{uniek leesbare} uitdrukkingen op. Daardoor
kunnen we met dezelfde methode---een inductieve definitie---betekenissen
aan die uitdrukkingen toekennen.

Het toekennen van betekenis aan uitdrukkingen behoort tot de semantiek.
Het centrale begrip in de semantiek van de propositielogica is
vervulling door !!a{valuation}. !!^a{valuation}~$\pAssign{v}$ kent de
waarheidswaarden $\True$ en $\False$ toe aan de !!{propositional
variable}s. Elke !!{valuation} bepaalt de waarheidswaarde
$\pValue{v}(!A)$ van elke !!{formula}~$!A$. !!^a{formula} wordt door
!!a{valuation}~$\pAssign{v}$ vervuld dan en slechts dan als
$\pValue{v}(!A) = \True$; dit noteren we als $\pSat{v}{!A}$. Deze
relatie kan ook door inductie naar de structuur van~$!A$ worden
gedefinieerd. Met de waarheidsfuncties van de logische connectieven
definiëren we bijvoorbeeld de vervulling van $!A \land !B$ in termen
van het al dan niet vervuld zijn van $!A$ en~$!B$.

Op grond van de vervullingsrelatie $\pSat{v}{!A}$ voor zinnen kunnen we
vervolgens de fundamentele semantische begrippen tautologie, semantisch
gevolg en vervulbaarheid definiëren. !!^a{formula} is een tautologie,
$\Entails !A$, als elke !!{valuation} haar vervult, dat wil zeggen als
$\pValue{v}(!A) = \True$ voor elke~$\pAssign{v}$. Een formule volgt uit
een verzameling !!{formula}s, $\Gamma \Entails !A$, als elke
!!{valuation} die alle !!{formula}s in~$\Gamma$ vervult, ook~$!A$
vervult. Een verzameling !!{formula}s is vervulbaar als er een
!!{valuation} is die alle !!{formula}s uit die verzameling tegelijk
vervult. Omdat !!{formula}s inductief zijn gedefinieerd en vervulling
op haar beurt door inductie naar de structuur van !!{formula}s is
gedefinieerd, kunnen we inductie gebruiken om eigenschappen van onze
semantiek te bewijzen en de gedefinieerde semantische begrippen met
elkaar in verband te brengen.
\end{document}
